\SourcePage{175}
\SGAchapter{XV}{Picard--Lefschetz 公式}{La formule de Picard-Lefschetz}
\begin{center}P. Deligne\end{center}
% 原讲目录由译本的节标题与页码重建。
\SourcePage{176}
\SGAsection{1}{普通二次奇点（典范形式）}
本节重新叙述第 VI 讲第 6 节的结果，采用“Hensel”的观点，而不采用“\Term{form}{形式}”的观点。我们给出\Term{ordinary-quadratic-singularity}{普通二次奇点}及其\Term{deformation}{形变}的\Term{canonical}{典范}\Term{form}{形式}；在\Term{etale-topology}{平展拓扑}下，它们\Term{local-scheme}{局部}都可化为这些\Term{form}{形式}。

这些结果很容易由第 VI 讲第 6 节、Tougeron \Term{form}{形式}的\Term{implicit-function-theorem}{隐函数定理}，以及 R. Elkik 关于\Term{isolated-singularity}{孤立奇点}的 Hensel \Term{versal-deformation}{万有形变}\Term{space}{空间}的\Term{theorem}{定理}\Term{pushout}{推出}。口头报告时，最后这个\Term{theorem}{定理}尚未问世；当时证明了一个较弱的\Term{theorem}{定理}，只涵盖\Term{complete-intersection}{完全交}的情形（这已足够我们使用），这个证明将不再发表。

\SGAsection{1.1}{回顾}
\textbf{1.1.1。} 设 $f:X\to S$ 是一个\Term{morphism}{态射}，$X$ \Term{local-scheme}{局部}可在 $S$ 上\Term{dimension}{维数}为 $n+p$ 的\Term{relative-affine-space}{相对仿射空间} $\mathbb A_S^{n+p}$ 中用 $p$ 个方程描述。\Term{jacobian-ideal}{Jacobi 理想} $J^n(X/S)$ 已在第 VI 讲第 5 节定义。若 $X$ 是由方程 $f_i=0$（$1\leq i\leq p$）定义的 $\mathbb A_S^{n+p}$ 的\Term{subscheme}{子概形}，它就是由 \Term{jacobian-matrix}{Jacobi 矩阵} $(\partial f_i/\partial x_j)$ 的 $p\times p$ 子式生成的\Term{ideal-sheaf}{理想层}。下面的\Term{theorem}{定理}是熟知的。它的\Term{form}{形式}类比的一个特例见 Bourbaki，\textit{Alg. Comm.} III，§4，第 6 小节，\Term{theorem}{定理} 2 的推论 3；$C^\infty$ 类比见 Tougeron，\textit{Idéaux de fonctions différentiables}（博士论文，1967 年 5 月，Rennes）；这里考虑的情形见 Artin，\textit{Algebraic approximation of structures over complete local rings}，\textit{Publ. Math. IHES} \textbf{36}（Zariski 专卷），\Term{lemma}{引理} 5.10。

\begin{SGAresult}{\Term{theorem}{定理}（Tougeron--Artin）}{1.1.2}
设 $S$ 是一个 Hensel \Term{local-ring}{局部环} $A$ 的\Term{spectrum}{谱}，$f:X\to S$ 是 1.1.1 所述的\Term{morphism}{态射}，$\mathfrak a$ 和 $I$ 是 $A$ 的两个真\Term{ideal}{理想}，$\delta'$ 是 $X$ 上包含于 $J^n(X/S)$ 的\Term{ideal-sheaf}{理想层}。假设 $I$ \Term{finite-type}{有限生成}。设 $s$ 是 $f$ 在 $S/I$ 上的一个\Term{section}{截面}：$s\in X(A/I)$。以 $s^{\#}\delta'^2$ 表示 $A/I$ 的\Term{ideal}{理想} $s^*\delta'^2$ 在 $A$ 中的\Term{inverse-image}{逆像}。如果
\[
 (s^{\#}\delta'^2)\mathfrak a\supset I,
\]
则 $f$ 有一个\Term{section}{截面} $s_1\in X(A)$，它模 $(s^*\delta')\mathfrak a$ 与 $s$ 同余。
\end{SGAresult}
\SourcePage{177}
\textbf{1.1.3。} 设 $S$ 是一个 Hensel \Term{local-ring}{局部环}的\Term{spectrum}{谱}，\Term{closed-point}{闭点}为 $s$。设 $X'_0$ 是 $k(s)$ 上的\Term{finite-type}{有限型}\Term{scheme}{概形}，在一个\Term{closed-point}{闭点} $x_0$ 之外\Term{smooth}{光滑}；令 $X_0$ 为 $X'_0$ 在 $x_0$ 处的\Term{henselization}{Hensel 化}。我们关心如下\Term{form}{形式}的系统：
\[
\begin{tikzcd}
X_0\arrow[r,hook]\arrow[d]&X\arrow[d]\\
s\arrow[r,hook]&T\arrow[r]&S .
\end{tikzcd}
\]
(a) $T$ 是\Term{local-scheme}{局部} Hensel 的，并且在 $S$ 上\Term{essentially-finitely-presented}{本质有限表示}（第 I 讲 0.0.5）；\Term{morphism}{态射} $T\to S$ 是\Term{local-scheme}{局部}的，并在\Term{residue-field}{剩余域}上诱导\Term{isomorphism}{同构}；

(b) $X$ 是\Term{local-scheme}{局部} Hensel 的，在 $T$ 上\Term{flat}{平坦}且\Term{essentially-finitely-presented}{本质有限表示}；

(c) $X/T$ 的\Term{special-fibre}{特殊纤维}配有一个到 $X_0$ 的 $k(s)$-\Term{isomorphism}{同构}。

这些系统称为 $X_0$ 的 \Term{henselian-deformation}{$S$-Hensel 形变}。从 $(X,T)$ 到 $(X',T')$ 的\Term{morphism}{态射}，是一个 $S$-\Term{scheme}{概形}的\Term{commutative-diagram}{交换图}
\[
\begin{tikzcd}X\arrow[r]\arrow[d]&X'\arrow[d]\\T\arrow[r]&T'\end{tikzcd}
\]
它在取\Term{special-fibre}{特殊纤维}后诱导 $X_0$ 的恒等。若\Term{deformation}{形变} $X/T$ 满足以下条件，就称其为\Term{versal-deformation}{万有的}（minimally versal）：

(a) 对任意\Term{morphism}{态射} $X''/T''\to X'/T'$，只要 $T''\to T'$ 是\Term{closed-immersion}{闭浸入}，映射
\[
\operatorname{Hom}(X'/T',X/T)\longrightarrow
\operatorname{Hom}(X''/T'',X/T)
\]
就是\Term{surjection}{满射}。特别地（取 $X''/T''=X_0/s$），$\operatorname{Hom}(X'/T',X/T)\ne\varnothing$。

(b) 对 $T'=\operatorname{Spec}(k(s)[\epsilon]/(\epsilon^2))$，“\Term{inverse-image}{逆像}”映射
\[
\operatorname{Hom}(T',T)\longrightarrow
\{X_0\text{ 在 }T'\text{ 上的形变的同构类}\}
\]
是\Term{bijection}{双射}。
\SourcePage{178}
\begin{SGAresult}{\Term{theorem}{定理}（R. Elkik）}{1.1.4}
$X_0$ 存在一个\Term{versal-deformation}{万有}的 \Term{henselian-deformation}{$S$-Hensel 形变}（在非唯一的\Term{isomorphism}{同构}意义下唯一）。若 $S$ 是\Term{noetherian}{Noether 的}，则 $X/T$ 为这样的\Term{deformation}{形变}的充要条件是：$X$ 在 $T$ 的\Term{completion}{完备化} $T^\wedge$ 上诱导 $X_0$ 的一个\Term{formal-versal-deformation}{形式万有 $S$-形变}（minimally formal versal deformation，第 VI 讲 1.9 及第 VI 讲第 4 节）。
\end{SGAresult}
这个精细\Term{theorem}{定理}的证明参见 R. Elkik（ENS 讨论班，1971/72）。

\SGAsection{1.2}{普通二次奇点}
设 $y$ 是\Term{characteristic}{特征}为 $p$ 的域 $k$ 上一个\Term{finite-type}{有限型}\Term{scheme}{概形} $Y$ 的\Term{closed-point}{闭点}，$n$ 为 $Y$ 在 $y$ 处的\Term{dimension}{维数}。
\begin{SGAresult}{定义}{1.2.1（= 第 VI 讲 6.6）}
(i) 若 $k$ \Term{algebraically-closed}{代数闭}，当 $Y$ 在 $y$ 处的\Term{local-ring}{局部环}的\Term{completion}{完备化} $\widehat{\mathscr O}_{Y,y}$ \Term{isomorphism}{同构}于 $k[[x_1,\ldots,x_{n+1}]]$ 对一个\Term{formal-power-series}{形式幂级数} $f$ 生成的\Term{ideal}{理想}的商，且 $f$ 可写成
\[
 f(\underline x)=Q(\underline x)=+\text{阶数}>2\text{ 的项},\tag{4.2.1}
\]
其中 $Q(\underline x)$ 是一个\Term{ordinary-quadratic-form}{普通二次型}时，就称 $y$ 为 $Y$ 的\Term{ordinary-quadratic-point}{普通二次点}。

(ii) 若 $\bar k$ 是 $k$ 的一个\Term{algebraic-closure}{代数闭包}，当 $\bar Y=Y\otimes_k\bar k$ 中位于 $y$ 之上的点均为 $\bar Y$ 的\Term{ordinary-quadratic-point}{普通二次点}时，就称 $y$ 为 $Y$ 的\Term{ordinary-quadratic-point}{普通二次点}。
\end{SGAresult}
\textbf{1.2.2。} 若在这个定义中将“\Term{ordinary-quadratic-form}{普通二次型}”换成“\Term{nondegenerate}{非退化}\Term{quadratic-form}{二次型}”，就得到\Term{nondegenerate-quadratic-point}{非退化二次点}的概念（第 VI 讲 6.0）。$y$ 是 $Y$ 的\Term{nondegenerate-quadratic-point}{非退化二次点}的充要条件是：它是\Term{ordinary-quadratic-point}{普通二次点}，并且 $p\ne2$ 或 $n$ 为奇数。
\SourcePage{179}
\begin{SGAresult}{例}{1.2.3}
设 $Q(\underline X)$ 是 $k$ 上一个 $n+1$ 元\Term{nondegenerate}{非退化}\Term{quadratic-form}{二次型}。设 $Y_0$ 为 $\mathbb A_k^{n+1}$ 中由 $Q=0$ 定义的\Term{quadratic-cone}{二次锥}。则原点（记为 $y_0$）是 $Y_0$ 的一个\Term{nondegenerate-quadratic-point}{非退化二次点}。
\end{SGAresult}
\begin{SGAresult}{例}{1.2.4}
设 $p=2$ 且 $n=2m$。取 $a\in k$，设 $Q$ 为\Term{polynomial}{多项式}
\[
 Q(\underline X)=(X_0^2-a)+\sum_{0<i\leq j\leq2m}a_{ij}X_iX_j.
\]
假设 $2m$ 元\Term{quadratic-form}{二次型} $\sum_{0<i\leq j}a_{ij}X_iX_j$ \Term{nondegenerate}{非退化}，以 $y_0$ 表示 $\mathbb A_k^{n+1}$ 中坐标为 $(\sqrt a,0,\ldots,0)$ 的\Term{closed-point}{闭点}。则 $y_0$ 是由方程 $Q=0$ 定义的 $\mathbb A_k^{n+1}$ 的\Term{subscheme}{子概形} $Y_0$ 的一个\Term{ordinary-quadratic-point}{普通二次点}。
\end{SGAresult}
\begin{SGAresult}{例}{1.2.5}
设 $k'$ 是 $k$ 的一个\Term{finite-separable-extension}{有限可分扩张}，$Y_0$ 是 1.2.3（\Term{lie-bracket}{括号}中的相应情形为 1.2.4）所考虑的 $k'$-\Term{scheme}{概形}。若将 $Y_0$ 看作 $k$-\Term{scheme}{概形}，则 $y_0$ 是 $Y_0$ 的一个\Term{nondegenerate-quadratic-point}{非退化二次点}（\Term{lie-bracket}{括号}中的相应情形为\Term{ordinary-quadratic-point}{普通二次点}）。$k$ 的\Term{group-extension}{扩张} $k'$ 可与 $k(y_0)$ 的最大\Term{separable-subextension}{可分子扩张}等同。
\end{SGAresult}
\begin{SGAresult}{\Term{theorem}{定理}}{1.2.6}
假设 $y$ 是 $k$-\Term{scheme}{概形} $Y$ 的一个\Term{ordinary-quadratic-point}{普通二次点}，$k'$ 为 $k(y)$ 的最大\Term{separable-subextension}{可分子扩张}。则存在一个 1.2.3 或 1.2.4 型的 $k'$-\Term{scheme}{概形} $Y_0\subset\mathbb A_{k'}^{n+1}$，以及 $Y$ 在 $y$ 处的\Term{henselization}{Hensel 化} $Y_{(y)}$ 与 $Y_0$ 在 $y_0$ 处的\Term{henselization}{Hensel 化} $Y_{0(y_0)}$ 之间的一个 $k$-\Term{isomorphism}{同构} $\varphi$。
\end{SGAresult}
首先注意，$y$ 在 $Y$ 中有一个\Term{etale-neighbourhood}{平展邻域} $Y'$，它是 $k'$-\Term{scheme}{概形}：
\[
\begin{tikzcd}&Y'\arrow[d]\\y\arrow[ur]\arrow[r]&Y .\end{tikzcd}
\]
事实上，\Term{scheme}{概形} $Y_{(y)}$ \Term{canonical}{典范}地是一个 $k'$-\Term{scheme}{概形}。
\SourcePage{180}
以 $(k',Y',y)$ 代替 $(k,Y,y)$，就归结到 $k=k'$ 的情形，即 $k(y)$ 在 $k$ 上\Term{radicial}{纯不可分}的情形。

$k(y)$-\Term{vector-space}{向量空间} $(\Omega^1_{Y/k})_y$ 的\Term{dimension}{维数}为 $n+1$：只须将标量由 $k(y)$ \Term{group-extension}{扩张}到其\Term{algebraic-closure}{代数闭包}，再在 $y$ 处\Term{completion}{完备化}后验证，此时结论显然。因此 $y$ 在 $Y$ 中的一个\Term{neighbourhood}{邻域}是一个\Term{dimension}{维数}为 $n+1$ 的\Term{smooth}{光滑} $k$-\Term{scheme}{概形} $Z$ 的\Term{subscheme}{子概形}。而且，适当缩小 $Y$ 和 $Z$ 为 $y$ 的\Term{neighbourhood}{邻域}后，定义 $Y$ 的 $\mathscr O_Z$ 的\Term{ideal}{理想}是\Term{principal-ideal}{主理想}，由一个方程 $f$ 定义（理由同上）。
\begin{SGAresult}{\Term{lemma}{引理}}{1.2.7}
假设 $y$ 是 $Y$ 的一个\Term{nondegenerate-quadratic-point}{非退化二次点}，且 $k(y)$ 在 $k$ 上\Term{radicial}{纯不可分}。则在 $y$ 附近，$J^n_{Y/k}$（1.1.1）是\Term{cyclic-module}{单生成的}，集中于 $y$，在 $k$ 上的\Term{rank}{秩}为 $1$。特别地，$k(y)=k$。
\end{SGAresult}
只须在过渡到 $k$ 的\Term{algebraic-closure}{代数闭包}并在 $y$ 处\Term{completion}{完备化}后证明第一条断言。此时直接计算即可看出：采用 1.2.1 的记号，由 \Term{nakayama-lemma}{Nakayama 引理}，$\partial f/\partial x_i$（从而 $f$ 和各 $\partial f/\partial x_i$）生成\Term{maximal-ideal}{极大理想}。

现在证明 $y$ \Term{nondegenerate}{非退化}时的 1.2.6。我们可以并且将假设 $k(y)=k$，$Y$ 在一个\Term{smooth}{光滑}\Term{scheme}{概形} $Z$ 中由一个方程 $f$ 定义，并存在\Term{etale}{平展}\Term{morphism}{态射} $x:Z\to\mathbb A_k^{n+1}$（一个“\Term{local-coordinate}{局部坐标}系 $x_i$”），将 $y$ 映至 $0$。于是存在一个 $n+1$ 元\Term{nondegenerate}{非退化}\Term{quadratic-form}{二次型} $Q$，使得在 $y$ 处
\[
 f=Q(x_i)+\text{阶数}>2\text{ 的项}.
\]
令 $Y_0$ 如 4.3 中那样由 $Q$ 定义。

在 $Y$ 上，$x_i$ 构成方程 $Q(x_i)=0$ 的一个“近似解”：若 $\mathfrak m$ 为定义 $y$ 的\Term{ideal}{理想}，则
\[
 Q(x_i)\in\mathfrak m^3\qquad(\text{在 }Y\text{ 上}).
\]
\SourcePage{181}
由于 $Q$ \Term{nondegenerate}{非退化}，$\mathbb A_k^{n+1}$ 中由各 $\partial Q/\partial x_i$ 生成的\Term{ideal}{理想}就是 $0$ 处的\Term{maximal-ideal}{极大理想}，而它在 $x=(x_i):Y\to\mathbb A_k^{n+1}$ 下的\Term{inverse-image}{逆像}为 $\mathfrak m$。由\Term{implicit-function-theorem}{隐函数定理} 1.1.2（取 $\mathfrak a=\mathfrak m$ 和 $\delta'=J$），在\Term{henselization}{Hensel 化} $Y_{(y)}$ 上存在函数 $x'_i$，模 $\mathfrak m^2$（即该处的 $s^*\delta'\mathfrak a$）与 $x_i$ 同余，并满足 $Q(x'_i)=0$。它们给出所求的 $Y_{(y)}$ 与 $Y_{0(y_0)}$ 之间的\Term{isomorphism}{同构}。
\begin{SGAresult}{\Term{lemma}{引理}}{1.2.8}
假设 $y$ 是一个\Term{degeneration}{退化}的\Term{ordinary-quadratic-point}{普通二次点}（$p=2$，$n$ 为偶数），且 $k(y)$ 在 $k$ 上\Term{radicial}{纯不可分}。则在 $y$ 附近，$T^1_{Y/k}$ 是\Term{cyclic-module}{单生成的}，集中于 $y$，在 $k$ 上的\Term{rank}{秩}为 $2$。特别地，或有 $k(y)=k$，或 $k(y)$ \Term{isomorphism}{同构}于 $k(\sqrt a)$，其中 $a\in k$ 且 $a\notin k^2$。
\end{SGAresult}
如 1.2.7 那样处理：采用 1.2.1 的记号，只须验证
\[
 \dim_k\left(k[[X_1,\ldots,X_{n+1}]]/(f,\partial f/\partial X_i)\right)=2.
\]
在适当的坐标系 $x_0,x_1,\ldots,x_{2m}$ 中，\Term{formal-power-series}{形式幂级数} $f$ 可写成
\[
 f=x_0^2+\sum_{i=1}^m x_i x_{i+m}+R\qquad(\operatorname{ord}R\geq3).
\]
令 $\mathfrak n'$ 为 $f$ 及各 $\partial f/\partial x_i$ 生成的\Term{ideal}{理想}，$\mathfrak n$ 为 $x_0^2$ 及各 $x_i$（$i\ne0$）生成的\Term{ideal}{理想}，$\mathfrak q$ 为\Term{maximal-ideal}{极大理想}。证明 $\mathfrak n=\mathfrak n'$。显然 $\mathfrak n'\subset\mathfrak n$；由 \Term{nakayama-lemma}{Nakayama 引理}，只须验证 $\mathfrak n'+\mathfrak n\mathfrak q=\mathfrak n$。模 $\mathfrak n\mathfrak q$ 有 $f\equiv x_0^2$；模 $\mathfrak n\mathfrak q+(x_0^2)$ 有
\[
 \partial f/\partial x_i\equiv x_{i+m},\qquad
 \partial f/\partial x_{i+m}\equiv x_i\quad(1\leq i\leq m),
\]
从而得到断言。

在 1.2.8 的假设下，若 $k=k(y)$，以 $a$ 表示 $k$ 中任意一个平方。
\SourcePage{182}
\begin{SGAresult}{\Term{lemma}{引理}}{1.2.9}
设 $V$ 是\Term{scheme}{概形} $S$ 上的一个\Term{vector-bundle}{向量丛}，$T$ 是 $\mathbb E_S(V)$ 的一个\Term{subscheme}{子概形}，在 $S$ 上\Term{locally-free}{局部自由}且\Term{rank}{秩}为 $2$。则 $T$ 包含在 $\mathbb E_S(V)$ 的唯一一条直线中。
\end{SGAresult}
这个问题在 $S$ 上是\Term{local-scheme}{局部}（fppf）的，因此可假设 $T/S$ 有一个\Term{section}{截面}。通过\Term{translation}{平移}，可假设 $T$ 包含 $0$ \Term{section}{截面}。设 $\mathfrak n$ 为 $T$ 中定义 $0$ 的\Term{ideal}{理想}：$\mathfrak n$ 在 $S$ 上\Term{locally-free}{局部自由}，\Term{rank}{秩}为 $1$。$\mathbb E_S(V)$ 上每个\Term{homogeneous}{齐次}\Term{linear-form}{线性形式} $v\in V$ 都定义 $\mathfrak n$ 的一个元素，从而得到 $V$ 到 $\mathfrak n$ 的一个\Term{epimorphism}{满态射}。其\Term{kernel}{核}定义所求直线。

\textbf{1.2.10。} 特别地，若 $A$ 是\Term{characteristic}{特征} $2$ 的域 $k$ 上\Term{affine-space}{仿射空间} $\mathbb A_k^{n+1}$ 的一个\Term{rank}{秩}为 $2$ 的\Term{radicial}{根性}\Term{subscheme}{子概形}，则存在 $a\in k$，使得在适当的\Term{affine}{仿射}坐标系中，$A$ 的方程为
\[
 x_0^2-a=0,\qquad x_i=0\quad(1\leq i\leq n).
\]
只须选择在 $A$ 确定的直线上为零的 $x_i$ 即可。

\textbf{1.2.11。} 证明 $y$ \Term{degeneration}{退化}时的 1.2.6（$p=2$，$n=2m$）。我们可以并且将假设 $k(y)=k(\sqrt a)$，$Y$ 在一个\Term{smooth}{光滑}\Term{scheme}{概形} $Z$ 中由一个方程 $f$ 定义，且存在 $y$ 处的\Term{local-coordinate}{局部坐标}系
\[
 x=(x_i)_{0\leq i\leq2m}:Z\longrightarrow\mathbb A_k^{2m+1},
\]
它将 $y$ 映至\Term{closed-point}{闭点} $(\sqrt a,0,\ldots,0)$，使 $J^n_{Y/k}$ 在 $y$ 附近\Term{isomorphism}{同构}于
$\mathscr O_Y/(x_0^2-a,x_1,\ldots,x_{2m})$。

于是，将 $f$ 换成适当的倍数后，有
\[
 f=(x_0^2-a)+\sum_{1\leq i\leq j\leq2m}a_{ij}x_i x_j+R,
\]
其中：

(a) $2m$ 元\Term{quadratic-form}{二次型} $\sum a_{ij}x_i x_j$ \Term{nondegenerate}{非退化}；

(b) 将\Term{scalar-extension}{标量扩张}到 $\bar k$ 后，$R$ 对 $x_0-\sqrt a$ 和各 $x_i$ 的阶数至少为 $3$。
\SourcePage{183}
如 1.2.4 那样，令 $Y_0$ 由
\[
 Q=x_0^2-a+\sum a_{ij}x_i x_j
\]
定义。在 $Y$ 上，仍可将 $x_i$ 看作方程 $Q(x_i)=0$ 的一个近似解。这一次，由各 $\partial Q/\partial X_i$ 生成的\Term{ideal}{理想}在 $x$ 下的\Term{inverse-image}{逆像} $\delta$ 满足 $\mathscr O_Y/\delta\simeq J^n_{Y/k}$，即在 $y$ 处
\[
 \delta=(x_0^2-a,(x_i)_{1\leq i}).
\]
设 $\mathfrak q$ 为定义 $y$ 的\Term{maximal-ideal}{极大理想}。为像\Term{nondegenerate}{非退化}情形那样应用\Term{implicit-function-theorem}{隐函数定理}而完成证明，须验证
\[
 Q(x_i)\equiv0\pmod{\delta^2\mathfrak q}\qquad\text{在 }Y\text{ 上}.\tag{1.2.11.1}
\]
设 $\bar{\mathfrak q}$ 为将\Term{scalar-extension}{标量扩张}到 $k$ 的\Term{algebraic-closure}{代数闭包}后定义 $y$ 的\Term{ideal}{理想}。可以验证
\[
 \delta^2\bar{\mathfrak q}\cap\mathscr O_{Y,y}=\delta^2\mathfrak q,
\]
所以只须将\Term{scalar-extension}{标量扩张}到 $\bar k$ 后验证 4.11.1：可假设 $a=0$。在 $y$ 处有
\[
 \delta^2\mathfrak q=
 \bigl(x_0^5,(x_0^3x_i)_{1\leq i},(x_0x_ix_j)_{1\leq i,j},
 (x_ix_jx_k)_{1\leq i,j,k}\bigr).
\]
由此可得 $\delta^2\mathfrak q\supset\mathfrak q^5$，以及
\[
\begin{aligned}
x_0^2f&\equiv x_0^4\pmod{\delta^2\mathfrak q},
&\text{故 }\delta^2\mathfrak q+(x_0^2f)&\supset\mathfrak q^4,\\
x_0f&\equiv x_0^3\pmod{\delta^2\mathfrak q+\mathfrak q^4},
&\text{故 }\delta^2\mathfrak q+(x_0f)&\supset\mathfrak q^3.
\end{aligned}
\]
\SourcePage{184}
最后有
\[
 Q(x_i)\equiv0\pmod{\delta^2\mathfrak q+(f)}\qquad\text{在 }Z\text{ 上},
\]
这就是 1.2.11.1。

\begin{SGAresult}{注记}{1.2.12}
设 $Q$ 是一个 $n+1$ 元、未必\Term{homogeneous}{齐次}且不恒为零的\Term{quadratic-form}{二次型}，假设由 $Q$ 定义的 $\mathbb A_k^{n+1}$ 的\Term{subscheme}{子概形} $Y$ 在某点 $y$ 有一个\Term{ordinary-quadratic-singularity}{普通二次奇点}。通过一个\Term{linear}{线性}（未必\Term{homogeneous}{齐次}）变量替换，就可将 $Y$ 化为 1.2.3 或 1.2.4 的\Term{form}{形式}。在\Term{nondegenerate}{非退化}情形，$k(y)=k$，只须将 $y$ 置于坐标原点。在\Term{degeneration}{退化}情形，只须使 $T^1_{Y/k}$ 的\Term{support}{支集}位于直线 $x_i=0$（$1\leq i\leq2m$）上（参见 1.2.9）。
\end{SGAresult}

\SGAsection{1.3}{二次奇点的模空间}
下面的命题由 R. Elkik 的\Term{theorem}{定理} 1.1.4 和留给读者的显式计算\Term{pushout}{推出}（计算实质上已包含在第 VI 讲第 6 节中）。
\begin{SGAresult}{命题}{1.3.1}
设 $S=\operatorname{Spec}(A)$ 是一个\Term{local-scheme}{局部} Hensel \Term{scheme}{概形}，\Term{closed-point}{闭点}为 $s$；设 $Y_0\subset\mathbb A_{k(s)}^{n+1}$ 为 1.2.3 或 1.2.4 型的\Term{scheme}{概形}。

(i) 取 1.2.3 型的 $Y_0$，由方程
\[
 Q(\underline X)=\sum_{i\leq j}\bar a_{ij}X_iX_j=0
 \qquad(\bar a_{ij}\in k(s))
\]
定义，将各 $\bar a_{ij}$ \Term{lifting}{提升}为 $A$ 中的元素 $a_{ij}$。则 $Y_0$ 在 $y_0$ 处\Term{singular-point}{奇点}的\Term{deformation}{形变}在 $S$ 上的 Hensel \Term{versal-deformation}{万有}\Term{base}{基底}\Term{scheme}{概形}为
\[
 T=\operatorname{Spec}(A\{b\}),
\]
\Term{versal-deformation}{万有形变}是 $T$-\Term{scheme}{概形} $Y\subset\mathbb A_T^{n+1}$ 在 $y_0$ 处的\Term{henselization}{Hensel 化}，其中 $Y$ 的方程为
\SourcePage{185}
\[
 \sum a_{ij}x_i x_j-b=0.
\]
(ii) 取 1.2.4 型的 $Y_0$，由方程
\[
 Q(\underline x)=(x_0^2-\bar a)+\sum\bar a_{ij}x_i x_j=0
 \qquad(\bar a,\bar a_{ij}\in k(s))
\]
定义，将 $\bar a_{ij}$ 和 $\bar a$ \Term{lifting}{提升}为 $A$ 中的元素 $a_{ij}$ 和 $a$。则 $Y_0$ 在 $y_0$ 处\Term{singular-point}{奇点}的\Term{deformation}{形变}在 $S$ 上的 Hensel \Term{versal-deformation}{万有}\Term{base}{基底}\Term{scheme}{概形}为
\[
 T=\operatorname{Spec}(A\{b,c\}),
\]
\Term{versal-deformation}{万有形变}是 $T$-\Term{scheme}{概形} $Y\subset\mathbb A_T^{n+1}$ 在 $y_0$ 处的\Term{henselization}{Hensel 化}，其中 $Y$ 的方程为
\[
 (x_0^2-a)+\sum a_{ij}x_i x_j+bx_0+c=0.
\]
\end{SGAresult}
下面的命题立即由 1.2.6 和 1.3.1 得到。
\begin{SGAresult}{推论}{1.3.2}
设 $S=\operatorname{Spec}(A)$ 是一个\Term{local-scheme}{局部} Hensel \Term{scheme}{概形}，\Term{closed-point}{闭点}为 $s$，$f:X\to S$ 是一个\Term{flat}{平坦}\Term{finitely-presented}{有限表示}\Term{morphism}{态射}。假设 $X_s$ 在其\Term{closed-point}{闭点} $x$ 处有一个\Term{ordinary-quadratic-singularity}{普通二次奇点}，$k(x)$ 在 $k(s)$ 上\Term{radicial}{纯不可分}，且 $X_s$ 的\Term{dimension}{维数}为 $n$。

(i) 若 $x$ \Term{nondegenerate}{非退化}，则存在一个系数在 $A$ 中的 $n+1$ 元\Term{nondegenerate}{非退化}\Term{quadratic-form}{二次型} $Q$，以及 $A$ 的\Term{maximal-ideal}{极大理想}中的 $b$，使得 $X$ 在 $x$ 处的\Term{henselization}{Hensel 化}这个 $S$-\Term{scheme}{概形}\Term{isomorphism}{同构}于 $\mathbb A_S^{n+1}$ 中方程（1.3.1(i)）$Q-b=0$ 所定义\Term{subscheme}{子概形}在原点的\Term{henselization}{Hensel 化}。

(ii) 若 $x$ \Term{degeneration}{退化}（$n=2m$，$\operatorname{char}k(s)=2$），则存在一个系数在 $A$ 中的 $n+1$ 元非\Term{homogeneous}{齐次}\Term{quadratic-form}{二次型}
\[
 Q(\underline X)=x_0+bx_0+c+\sum_{1\leq i\leq j\leq2m}a_{ij}x_i x_j,\tag{4.14.2}
\]
\SourcePage{186}
其中 $b$ 属于\Term{maximal-ideal}{极大理想}，$2m$ 元\Term{quadratic-form}{二次型} $\sum a_{ij}x_i x_j$ \Term{nondegenerate}{非退化}，并使得 $X$ 在 $x$ 处的\Term{henselization}{Hensel 化}这个 $S$-\Term{scheme}{概形}\Term{isomorphism}{同构}于 $\mathbb A_S^{n+1}$ 中方程 $Q=0$ 所定义\Term{subscheme}{子概形}在 $(\sqrt c,0,\ldots,0)$ 处的\Term{henselization}{Hensel 化}。
\end{SGAresult}
\begin{SGAresult}{注记}{1.3.3}
在 (i) 的情形，若 $X_s$ 在 $x$ 处的\Term{henselization}{Hensel 化}通过一个\Term{isomorphism}{同构} $\varphi$ \Term{isomorphism}{同构}于 $X_0\subset\mathbb A_k^{n+1}$ 在 $0$ 处的\Term{henselization}{Hensel 化}，其中 $X_0$ 的方程为 $\sum\bar a_{ij}x_i x_j=0$，且各 $a_{ij}$ \Term{lifting}{提升}各 $\bar a_{ij}$，则可取 $Q=\sum a_{ij}x_i x_j$。

在 (ii) 的情形，若 $X_s$ 在 $x$ 处的\Term{henselization}{Hensel 化}通过一个\Term{isomorphism}{同构} $\varphi$ \Term{isomorphism}{同构}于 $X_0\subset\mathbb A_k^{n+1}$ 在 $(\sqrt{\bar a},0,\ldots,0)$ 处的\Term{henselization}{Hensel 化}，其中 $X_0$ 的方程为 $(x_0^2+\bar a)+\sum\bar a_{ij}x_i x_j$，且各 $a_{ij}$ \Term{lifting}{提升}各 $\bar a_{ij}$，则可取
\[
 Q=x_0^2+bx_0+c+\sum a_{ij}x_i x_j,
 \qquad c\equiv\bar a\pmod{\text{极大理想}}.
\]
两种情形下，1.3.2 所保证的\Term{isomorphism}{同构}都可取为延拓 $\varphi$ 的\Term{isomorphism}{同构}。
\end{SGAresult}
\begin{SGAresult}{推论}{1.3.4}
设 $f:X\to S$ 如 1.3.2 所述。则 $x$ 在 $X$ 中有一个\Term{neighbourhood}{邻域} $U$，使得 $U$ 中 $f$ 不\Term{smooth}{光滑}的点，都是其所在\Term{fibre}{纤维}的\Term{ordinary-quadratic-point}{普通二次点}。
\end{SGAresult}
\SourcePage{187}
\SGAsection{2}{一个标准普通二次情形中消失圈的计算}
\SGAsection{2.1}{预备：锥的上同调}
\textbf{2.1.1。} 设 $Y\subset\mathbb P^r$ 是\Term{algebraically-closed}{代数闭}域 $k$ 上的一个\Term{projective-variety}{射影簇}。采用以下记号。

$X\subset\mathbb A^{r+1}$ 为以 $0$ 为顶点的\Term{affine}{仿射}\Term{projecting-cone}{投影锥}；

$X_{(0)}$ 为 $X$ 在 $0$ 处的\Term{henselization}{Hensel 化}；

$X^*=X-\{0\}$，$X^*_{(0)}=X_{(0)}-\{0\}$；

$X_1\subset\mathbb P^{r+1}$ 为 $Y$ 的\Term{projective}{射影}\Term{projecting-cone}{投影锥}；有 $X=X_1-Y$；

$\widetilde X$、$\widetilde X_{(0)}$、$\widetilde X_1$ 分别为 $X$、$X_{(0)}$、$X_1$ 在 $0$ 处的\Term{blow-up}{爆破}；

$Y_0$ 为 $0$ 在 $\widetilde X$、$\widetilde X_{(0)}$ 或 $\widetilde X_1$ 中的\Term{inverse-image}{逆像}；

$h:\widetilde X_1\to Y$ 为自然\Term{projection}{投影}；

$y_0,y_\infty:Y\to\widetilde X_1$ 为 $h$ 的\Term{section}{截面}，其像分别为 $Y_0$ 和 $Y$；

$F$ 为一个\Term{torsion-abelian-group}{挠阿贝尔群}，其\Term{torsion}{挠性}与 $k$ 的\Term{characteristic}{特征}互素；所考虑的所有\Term{cohomology}{上同调}群均以 $F$ 为系数，记号中略去这个系数。

本小节的目标是联系 $X^*_{(0)}$ 和 $Y$ 的\Term{cohomology}{上同调}。
\begin{SGAresult}{命题}{2.1.2}
有
\[
\text{(i)}\quad H^i(X)\xrightarrow{\sim}H^i(\{0\})=
\begin{cases}F&i=0,\\0&i>0;\end{cases}
\qquad
\text{(ii)}\quad H^i_{\{0\}}(X)\xrightarrow{\sim}H_c^i(X).
\]
\end{SGAresult}
这个命题可由一个\Term{homotopy}{同伦}论证证明；我们须先将该论证\Term{formalization}{形式化}。

设 $k$ 是一个\Term{algebraically-closed}{代数闭}域，$U$ 和 $V$ 是两个 $k$-\Term{scheme}{概形}，$\Lambda$ 是一个\Term{torsion}{挠性}与 $k$ 的\Term{characteristic}{特征}互素的\Term{torsion-ring}{挠环}，$K\in\operatorname{Ob}D^+(U,\Lambda)$，$L\in\operatorname{Ob}D^+(V,\Lambda)$。
\SourcePage{188}
从 $(U,K)$ 到 $(V,L)$ 的一个“\Term{morphism}{态射}” $f$，是由 $f:U\to V$ 和 $\varphi:f^*L\to K$ 组成的一对。这样的\Term{morphism}{态射}在\Term{cohomology}{上同调}上诱导
\[
 f^*:H^*(V,L)\longrightarrow H^*(U,K).
\]
若存在 $k$ 上的一个\Term{connected}{连通}\Term{finite-type}{有限型}\Term{scheme}{概形} $T$、$T$ 的两个点 $0$ 和 $1$，以及一个从 $(U\times T,\operatorname{pr}_1^*K)$ 到 $(V,L)$ 的“\Term{morphism}{态射}” $H$，取 $0$ 和 $1$ 处的\Term{fibre}{纤维}后分别诱导 $f_0$ 和 $f_1$，就说这两个“\Term{morphism}{态射}” $f_0$ 与 $f_1$ \Term{homotopy}{同伦}。
\begin{SGAresult}{\Term{lemma}{引理}}{2.1.3}
若 $f_0$ 与 $f_1$ \Term{homotopy}{同伦}，则 $f_0^*=f_1^*$。
\end{SGAresult}
已知存在 $T$ 中的一列点 $0=x_0,x_1,\ldots,x_n=1$、\Term{smooth}{光滑}\Term{connected}{连通}\Term{curve}{曲线} $\Gamma_i$ 和\Term{morphism}{态射} $a_i:\Gamma_i\to T$，使得 $x_i$ 和 $x_{i+1}$ 都在 $a_i$ 的像中。为构造各 $\Gamma_i$，对 $T$ 中一个包含 $0$ 和 $1$ 的一维\Term{connected}{连通}\Term{subscheme}{子概形}作\Term{normalization}{正规化}即可。因此可归结到 $T$ 为\Term{smooth}{光滑}\Term{curve}{曲线}的情形。考虑
\[
\begin{tikzcd}
U\times T\arrow[r,"\operatorname{pr}_1"]\arrow[d,"\operatorname{pr}_2"']&U\arrow[d,"f"]\\
T\arrow[r,"t"]&k .
\end{tikzcd}
\]
对 $t$ 应用\Term{smooth-base-change}{光滑基变换定理}，得到
\[
 t^*Rf_*K\xrightarrow{\sim}R\operatorname{pr}_{2*}(\operatorname{pr}_1^*K).
\]
于是映射 $f_i^*$（$i=0,1$）\Term{resolution}{分解}为
\[
\begin{tikzcd}[column sep=small]
H^n(V,L)\arrow[r,"H^*"]&H^n(U\times T,\operatorname{pr}_1^*K)\arrow[r]&
H^0(T,R^n\operatorname{pr}_{2*}(\operatorname{pr}_1^*K))\arrow[dd,"\text{在 }i\text{ 处的纤维}"]\\
&H^0(T,t^*H^n(U,K))\arrow[ur,"\sim"]\arrow[dr,"\sim"']&\\
&&H^n(U,K)
\end{tikzcd}
\]
其中 $H^0(T,t^*H^n(U,K))$ 与右上角通过上述\Term{isomorphism}{同构}等同，与 $H^n(U,K)$ 通过常值\Term{section}{截面}的取值等同，故 $f_i^*$ 与 $i$ 无关。
\SourcePage{189}
\textbf{2.1.3。} 证明 2.1.2。借助\Term{homothety}{位似}，$X$ 的\Term{identity-map}{恒等映射}与 $X\to X$ 中值为 $0$ 的常值映射\Term{homotopy}{同伦}；由此得到 (i)。

用比例趋于无穷大的\Term{homothety}{位似}，可见 $Y$ 是 $X_1^*$ 的一个\Term{deformation-retract}{形变收缩核}；由\Term{five-lemma}{五引理}和下图即得 (ii)：
\[
\begin{tikzcd}[column sep=small]
\cdots\arrow[r,"\partial"]&H^i_{\{0\}}(X)\arrow[r]\arrow[d]&H^i(X_1)\arrow[r]\arrow[d,equal]&H^i(X_1^*)\arrow[r,"\partial"]\arrow[d,"\sim"]&\cdots\\
\cdots\arrow[r,"\partial"]&H_c^i(X)\arrow[r]&H^i(X_1)\arrow[r]&H^i(Y)\arrow[r,"\partial"]&\cdots .
\end{tikzcd}
\]
\begin{SGAresult}{推论}{2.1.4}
$H^i(X^*)\xrightarrow{\sim}H^i(X^*_{(0)})$。
\end{SGAresult}
对下图应用\Term{five-lemma}{五引理}：
\[
\begin{tikzcd}[column sep=small]
\cdots\arrow[r]&H^i_{\{0\}}(X)\arrow[r]\arrow[d,"\sim"]&H^i(X)\arrow[r]\arrow[d,"{\text{2.1.2(i)}}","\sim"']&H^i(X^*)\arrow[r]\arrow[d]&\cdots\\
\cdots\arrow[r]&H^i_{\{0\}}(X_{(0)})\arrow[r]&H^i(X_{(0)})\arrow[r]&H^i(X^*_{(0)})\arrow[r]&\cdots .
\end{tikzcd}
\]
\textbf{2.1.5。} 利用 $h|\widetilde X_1-Y_0$ 及 $h|\widetilde X_1-Y$ 的 Leray \Term{spectral-sequence}{谱序列}，可得\Term{restriction-morphism}{限制态射}
\[
 H^i(\widetilde X_1-Y_0)\xrightarrow{\sim}H^i(Y),\qquad
 H^i(\widetilde X)\xrightarrow{\sim}H^i(Y_0)=H^i(Y)
\]
均为\Term{isomorphism}{同构}。通过这些\Term{isomorphism}{同构}，对 $(\widetilde X_1-Y_0,Y)$、$(\widetilde X_1-Y,Y_0)$ 的\Term{cohomology}{上同调}\Term{long-exact-sequence}{长正合列}成为下图的两行：
\SourcePage{190}
\[
\begin{tikzcd}[column sep=small]
H^{i-1}(X^*)\arrow[r]\arrow[d,equal]&H^{i-2}(Y)(-1)\arrow[r,"\eta"]\arrow[d,equal]&H^i(Y)\arrow[r]\arrow[d,"-1"]&H^i(X^*)\arrow[r]\arrow[d,equal]&\cdots\\
H^{i-1}(X^*)\arrow[r]&H^{i-2}(Y)(-1)\arrow[r,"-\eta"]&H^i(Y)\arrow[r]&H^i(X^*)\arrow[r]&\cdots .
\end{tikzcd}\tag{2.1.5.1}
\]
\begin{SGAresult}{\Term{lemma}{引理}}{2.1.6}
2.1.5.1 中记为 $\eta$ 和 $-\eta$ 的\Term{arrow}{箭头}是与一个\Term{hyperplane-section}{超平面截面}的\Term{cohomology-class}{上同调类}的\Term{cup-product}{杯积}。图 2.1.5.1 \Term{commutative-diagram}{交换}。
\end{SGAresult}
第一条断言来自：$\eta$（相应地 $-\eta$）是 $\widetilde X_1-Y_0$（相应地 $\widetilde X_1-Y$）上\Term{invertible-sheaf}{可逆层} $\mathscr O(Y)$（相应地 $\mathscr O(Y_0)$）的\Term{cohomology-class}{上同调类}在 $Y$（相应地 $Y_0\simeq Y$）上的\Term{restriction}{限制}。

利用 $h$ 的 Leray \Term{spectral-sequence}{谱序列}：(2.1.5.1) 的两行是对下列\Term{distinguished-triangle}{区别三角}应用\Term{functor}{函子} $H(Y,\ )$ 得到的：
\[
 y_{\infty*}Ry_\infty^!F\longrightarrow R(h|\widetilde X_1-Y_0)_*F
 \xrightarrow{\text{\textcircled{1}}}R(h|X^*)_*F\xrightarrow{\text{\textcircled{2}}},
\]
以及
\[
 y_{0*}Ry_0^!F\longrightarrow R(h|\widetilde X_1-Y)_*F
 \longrightarrow R(h|X^*)_*F\longrightarrow.
\]
这两个三角各顶点\Term{complex}{复形}的非零\Term{cohomology-sheaf}{上同调层}的\Term{degree}{次数}分别为 $(2)$、$(0)$、$(0\text{ 和 }1)$，所以\Term{arrow}{箭头} \textcircled{1} 和 \textcircled{2} 由它们在\Term{cohomology-sheaf}{上同调层}上的\Term{action}{作用}确定。因此归结到 $Y$ 为一点的情形；这一情形留给读者。

\textbf{2.1.7。} 对 $(\widetilde X_{(0)},Y_0)$ 的\Term{long-exact-sequence}{长正合列}是 (2.1.5.1) 第二行的\Term{local-scheme}{局部}类比。有
\[
 H^i(\widetilde X_{(0)})\xrightarrow{\sim}H^i(Y_0)
\]
（对\Term{proper}{固有}\Term{morphism}{态射} $\widetilde X_{(0)}\to X_{(0)}$ 应用\Term{proper-base-change}{固有基变换定理}）；因而该\Term{long-exact-sequence}{长正合列}可写成
\SourcePage{191}
\[
\begin{tikzcd}[column sep=small]
H^{i-1}(X^*_{(0)})\arrow[r]&H^{i-2}(Y_0)(-1)\arrow[r,"-\eta"]\arrow[d,equal]&H^i(Y_0)\arrow[r]\arrow[d,equal]&H^i(X^*_{(0)})\\
H^{i-1}(X^*)\arrow[r]\arrow[u,"\sim"]&H^{i-2}(Y)(-1)\arrow[r,"-\eta"]&H^i(Y)\arrow[r]&H^i(X^*)\arrow[u,"\sim"] .
\end{tikzcd}\tag{2.1.7.1}
\]
这个\Term{commutative-diagram}{交换图}的第二行就是 (2.1.5.1) 的第二行。

下面的技术\Term{lemma}{引理}将在 2.2.7 中使用。
\begin{SGAresult}{\Term{lemma}{引理}}{2.7.8}
下图\Term{anticommutative}{反交换}：
\[
\begin{tikzcd}[column sep=small]
H^{n-1}(Y_0)\arrow[r,"\sim",leftarrow]\arrow[d,equal]&H^{n-1}(\widetilde X)\arrow[r]&H^{n-1}(X^*)\arrow[r]&H^n_{\{x\}}(X)\arrow[r]&H_c^n(X_s)\arrow[d,equal]\\
H^{n-1}(Y)\arrow[rrrr,"\partial"']&&&&H_c^n(X_s).
\end{tikzcd}
\]
\end{SGAresult}
由 2.1.7，\Term{equivalence}{等价}地，只须证明下图\Term{commutative-diagram}{交换}：
\[
\begin{tikzcd}
H^{n-1}(X_1-\{0\})\arrow[r,"\partial"]\arrow[d,"\sim"]&H^n_{\{0\}}(X_1)=H^n_{\{0\}}(X)\arrow[d]\\
H^{n-1}(Y)\arrow[r,"\partial"]&H_c^n(X).
\end{tikzcd}
\]
这只是\Term{generic}{一般}的\Term{form}{形式}论证。

\SGAsection{2.2}{消失圈}
\textbf{2.2.1。} 设 $S$ 是 Hensel \Term{trait}{离散赋值环谱}，$s,\eta,\bar s,\bar\eta$ 如第 XIII 讲 0.2.5 所述；$\Lambda=\mathbb Z/k$，其中 $k$ 与 $S$ 的\Term{residue-characteristic}{剩余特征}互素；$\mathbb A_S^{n+1}$ 为 $S$ 上标准的 $n+1$ 维\Term{affine-space}{仿射空间}，$X$ 为由一个模\Term{uniformizer}{一致化参数}不为零的二次方程 $Q=0$ 定义的 $\mathbb A_S^{n+1}$ 的\Term{closed-subscheme}{闭子概形}，其中
\SourcePage{192}
\[
 Q(\underline x)=\sum_{i\leq j}a_{ij}x_i x_j+\sum b_i x_i+c.
\]
若 $\mathbb A_S^{n+1}$ 是 $\mathbb P_S^{n+1}$ 中无穷远\Term{hyperplane}{超平面} $H$ 的补，则 $X$ 由 $\mathbb P_S^{n+1}$ 中\Term{homogeneous}{齐次}方程
\[
 \sum_{i\leq j}a_{ij}x_i x_j+\sum b_i x_i z+cz^2=0\tag{2.2.1.1}
\]
定义的\Term{quadric}{二次超曲面} $X_1$ 去掉 $Y=X_1\cap H$ 得到。

作如下假设。

(a) $Y$ 是 $S$ 上的\Term{smooth}{光滑}\Term{quadric}{二次超曲面}，即\Term{quadratic-form}{二次型} $\sum_{i\leq j}a_{ij}x_i x_j$ 是\Term{ordinary-quadratic-form}{普通的}；

(b) \Term{quadric}{二次超曲面} $X'_s$ 是\Term{singular}{奇异的}，即 $X_{\bar s}$ 为\Term{quadratic-cone}{二次锥}。

以 $x_0$ 表示 $X_s$ 的\Term{singular-point}{奇点}，以 $f$ 表示 $X$ 到 $S$ 的\Term{projection}{投影}。
\begin{SGAresult}{注记}{2.2.2}
若 $x_0$ 是\Term{rational-point}{有理点}，可通过\Term{translation}{平移}将它置于坐标原点。若 $x_0=0$，则各 $b_i$ 和 $c$ 均被\Term{uniformizer}{一致化参数}整除。

$X_s$ 中由方程 $\partial Q/\partial x_i$ 定义的\Term{subscheme}{子概形} $A$ 集中于 $x_0$。除 $k(s)$ 的\Term{characteristic}{特征}为 $2$ 且 $n+1$ 为偶数这个例外情形外，这个\Term{scheme}{概形}在 $k(s)$ 上的\Term{degree}{次数}为 $1$，因而 $x_0$ 是\Term{rational-point}{有理点}。在例外情形，它在 $k(s)$ 上是\Term{ramified}{分歧的}且\Term{rank}{秩}为 $2$，故 $k(x_0)$ 或为 $k(s)$，或为 $k(s)$ 的一个二次不\Term{separable-extension}{可分扩张}。关于 $A$ 的这些断言，过渡到 $k(s)$ 的\Term{algebraic-closure}{代数闭包} $k(\bar s)$ 后最容易验证。
\end{SGAresult}
\begin{SGAresult}{命题}{2.2.3}
在 2.2.1 的假设下，\Term{morphism}{态射}
\SourcePage{193}
\[
 H^i(X_{\bar\eta},\Lambda)\xrightarrow{\sim}
 \mathbb H^i(X_{\bar s},R\Psi_{\bar\eta}(\Lambda))\xrightarrow{\sim}
 R^i\Psi_{\bar\eta}(\Lambda)_{x_0},
\]
\[
 H_c^i(X_{\bar\eta},\Lambda)\xleftarrow{\sim}
 \mathbb H_c^i(X_{\bar s},R\Psi_{\bar\eta}(\Lambda))\xleftarrow{\sim}
 H^i_{\{x_0\}}(X_{\bar s},R\Psi_{\bar\eta}(\Lambda))
\]
均为\Term{isomorphism}{同构}。
\end{SGAresult}
左侧\Term{arrow}{箭头}为\Term{isomorphism}{同构}由 2.1.8.6 和 2.1.10.5 得到；右侧\Term{arrow}{箭头}为\Term{isomorphism}{同构}是 2.1.2 的推论。事实上，有\Term{distinguished-triangle}{区别三角}
\[
 (\Lambda\text{ 在 }X_{\bar s}\text{ 上})[0]\longrightarrow
 R\Psi_{\bar\eta}(\Lambda)\longrightarrow R\Phi_{\bar\eta}(\Lambda),
\]
而 2.1.2 可用于 $(\Lambda\text{ 在 }X_{\bar s}\text{ 上})$，同时 $R\Phi_{\bar\eta}(\Lambda)$ 的\Term{support}{支集}在 $x_0$ 中。
\begin{SGAresult}{推论}{2.2.4}
若\Term{geometric-generic-fibre}{几何一般纤维} $X_{\bar\eta}$ 是\Term{singular}{奇异的}，即仍为\Term{quadratic-cone}{二次锥}，则所有\Term{sheaf}{层} $R^i\Phi_{\bar\eta}(\Lambda)$ 都为零。
\end{SGAresult}
由 2.1.2，$\Lambda=H^0(X_{\bar s},\Lambda)\xrightarrow{\sim}H^0(X_{\bar\eta},\Lambda)$，且 $0=H^i(X_{\bar s},\Lambda)=H^i(X_{\bar\eta},\Lambda)$。于是应用第 XIII 讲的\Term{long-exact-sequence}{长正合列} 2.1.8.9（由第 XIII 讲 2.1.9 可知此列适用）。

\textbf{2.2.5。} 现在假设

$(*)$ \Term{generic-fibre}{一般纤维} $X_\eta$ \Term{smooth}{光滑}。

为简便起见，再假设 $S$ 为\Term{strictly-henselian}{严格 Hensel 的}。利用 2.2.3 和第 XII 讲 3.7，容易计算\Term{vanishing-cycles}{消失圈}\Term{sheaf}{层}。这里考虑的以 $\Lambda$ 为系数的\Term{cohomology}{上同调}，由该处的 $\ell$-进\Term{cohomology}{上同调}通过\Term{versal}{万有}系数公式得到（参见 2.1.13）。结果如下。

\textbf{A. 不带\Term{support}{支集}的\Term{cohomology}{上同调}。}
有 $R^i\Psi_{\bar\eta}(\Lambda)=0$，若 $i\ne0,n$。
\SourcePage{194}
1) 若 $n\ne0$：有 $\Psi_{\bar\eta}(\Lambda)=\Lambda$。此外，$R^n\Psi_{\bar\eta}(\Lambda)$（非\Term{canonical}{典范}地）\Term{isomorphism}{同构}于 $x_0$ 上的\Term{sheaf}{层} $\Lambda$ 在 $X_s$ 上的\Term{extension-by-zero}{零延拓}。

2) 任意 $n$：可以验证，即使 $n=0$，当 $i\ne n$ 时仍有 $\Phi^i_{\bar\eta}(\Lambda)=0$，而 $\Phi^n_{\bar\eta}(\Lambda)$（非\Term{canonical}{典范}地）\Term{isomorphism}{同构}于 $x_0$ 上的\Term{sheaf}{层} $\Lambda$ 在 $X_s$ 上的\Term{extension-by-zero}{零延拓}。

\textbf{B. \Term{support}{支集}在 $x_0$ 中的\Term{cohomology}{上同调}。}
当 $i\ne n,2n$ 时，有 $H^i_{\{x_0\}}(X_s,R\Psi_{\bar\eta}(\Lambda))=0$。此外有\Term{trace-morphism}{迹态射}
\[
 \operatorname{tr}:H^{2n}_{\{x_0\}}(X_s,R\Psi_{\bar\eta}(\Lambda(n)))\longrightarrow\Lambda.\tag{2.2.5.1}
\]
若 $n\ne0$：\Term{morphism}{态射} (2.2.5.1) 是一个\Term{isomorphism}{同构}。此外，$H^n_{\{x_0\}}(X_s,R\Psi_{\bar\eta}(\Lambda(n)))$（非\Term{canonical}{典范}地）\Term{isomorphism}{同构}于 $\Lambda$。

\textbf{C. \Term{duality}{对偶性}。} 对 $a\in R^n\Psi_{\bar\eta}(\Lambda)_{x_0}$ 和 $b\in H^n_{\{x_0\}}(X_s,R\Psi_{\bar\eta}(\Lambda(n)))$，令 $(a,b)=\operatorname{Tr}(a\wedge b)$，其中 $\operatorname{Tr}$ 是第 XIII 讲 2.2.6 的\Term{morphism}{态射}
\[
 \operatorname{Tr}:H^{2n}_{\{x_0\}}(X_s,R\Psi_{\bar\eta}(\Lambda(n)))\longrightarrow\Lambda.
\]
由 $X_{\bar\eta}$ 上的 Poincaré \Term{duality}{对偶性}及 2.2.3 的\Term{isomorphism}{同构}，$(a,b)$ 使\Term{free-module}{自由} $\Lambda$-\Term{module}{模} $R^n\Psi_{\bar\eta}(\Lambda)_{x_0}$ 与 $H^n_{\{x_0\}}(X_s,R\Psi_{\bar\eta}(\Lambda(n)))$ 构成\Term{perfect-duality}{完美对偶}。

\textbf{D. $n=2m>0$。}
由第 XII 讲 3.7，群 $H^n_{\{x_0\}}(X_s,R\Psi_{\bar\eta}(\Lambda(m)))$ 和 $R^n\Psi_{\bar\eta}(\Lambda(m))_{x_0}$ 具有精确到\Term{signature}{符号}的自然\Term{generator}{生成元}，分别记为 $\delta$ 和 $\delta'$。设 $\varphi$ 为前一个群到后一个群的自然映射。可规范 $\delta$ 和 $\delta'$ 的\Term{signature}{符号}，使得
\[
 (\delta',\delta)=1;\tag{2.2.5.2}
\]
于是有
\SourcePage{195}
\[
 \varphi(\delta)=(-1)^m\,2\delta'.\tag{2.2.5.3}
\]
设 $Z$ 为\Term{quadratic-form}{二次型} (2.2.1.1) 的\Term{clifford-algebra}{Clifford 代数}的偶部的\Term{centre-algebra}{中心}，它是 $k(\eta)$ 的\Term{separable-quadratic-extension}{可分二次扩张}。设 $\epsilon$ 为\Term{inertia-group}{惯性群} $I$ 上对应的阶为 $1$ 或 $2$ 的\Term{character}{特征标}。对 $\sigma\in I$，有
\[
 \sigma\delta=\epsilon(\sigma)\delta,\tag{2.2.5.4}
\]
\[
 \sigma\delta'=\epsilon(\sigma)\delta'.\tag{2.2.5.5}
\]
在以 $\mathbb Z/2k$ 为系数的\Term{cohomology}{上同调}中，公式 (2.2.5.2)、(2.2.5.3)、(2.2.5.4) 蕴涵\Term{variation}{变差}为
\[
 \operatorname{Var}(\sigma)(a)=\frac{\epsilon(\sigma)-1}{2}(-1)^m(a\delta)\delta.\tag{2.2.5.6}
\]
\textbf{E. $n=2m+1$——\Term{vanishing-cycles}{消失圈}。}
由第 XII 讲 3.7，群 $H^n_{\{x_0\}}(X_s,R\Psi_{\bar\eta}(\Lambda(m)))$ 与 $R^n\Psi_{\bar\eta}(\Lambda(m+1))_{x_0}$ 具有精确到\Term{signature}{符号}的自然\Term{generator}{生成元}，分别记为 $\delta$ 和 $\delta'$。沿用 D 的记号，可规范 $\delta$ 和 $\delta'$ 的\Term{signature}{符号}，使得
\[
 (\delta',\delta)=1;\tag{2.2.5.7}
\]
此时有
\[
 \varphi(\delta)=0.\tag{2.2.5.8}
\]
\textbf{F. $n=2m+1$——\Term{variation}{变差}。}
\Term{inertia-group}{惯性群} $I$ 在 $X_{\bar\eta}$ 的\Term{cohomology}{上同调}上的\Term{action}{作用}平凡。这由第 XII 讲 3.7 以及 $I$ 在 $Y_{\bar\eta}$ 的\Term{cohomology}{上同调}上\Term{action}{作用}平凡得到；这里 $Y$ 在 $S$ 上\Term{proper}{固有}且\Term{smooth}{光滑}。

对\Term{inertia-group}{惯性群}中的 $\sigma$，必有
\[
 D(\sigma)(\delta')=\lambda(\sigma)\delta,\qquad\lambda(\sigma)\in\Lambda(1).\tag{2.2.5.9}
\]
可\Term{additive}{加性}公式（第 XIII 讲 1.4.3.4）在这里写成
\[
 \lambda(\sigma\tau)=\lambda(\sigma)+\lambda(\tau).
\]
令 $\epsilon:I\to\Lambda(1)=\mu_k$ 为满足下式的\Term{homomorphism}{同态}：
\SourcePage{196}
\[
 \sigma(\sqrt[k]{t})=\epsilon(\sigma)\sqrt[k]{t},
\]
其中 $t$ 为一个\Term{uniformizer}{一致化参数}。由\Term{tame-inertia-group}{温和惯性群}的结构，$\lambda$ 必为 $\epsilon$ 的一个倍数。公式 (2.2.5.9) 可重写为
\[
 \operatorname{Var}(\sigma)(a)=\lambda_X\epsilon(\sigma)(a\delta)\delta,\tag{2.2.5.10}
\]
其中 $\lambda_X\in\Lambda$ 适当选择（依赖于 $X/S$）。稍后将确定 $\lambda_X$。

\textbf{2.2.6。} 现在对 (D) 和 (E) 中的\Term{cycle}{圈} $\delta$ 和 $\delta'$ 给出一个在\Term{etale-topology}{平展拓扑}下为\Term{local-character}{局部性}的描述。

$n$ 为偶数的情形很容易：若 $2\nmid k$（$\Lambda=\mathbb Z/k$），则 $\pm\delta$ 由 (2.2.5.3)、(2.2.5.4) 刻画，即由
\[
 (\delta,\delta)=(-1)^m\,2\tag{2.2.6.1}
\]
刻画。\Term{generic}{一般}地，$\pm\delta$ 是以 $\mathbb Z/2^ak$（$a$ \Term{large-relative-to-cardinal}{足够大}）为值且满足 (2.2.6.1) 的任意\Term{cycle}{圈} $\widetilde\delta$ 的像。

现在假设 $n$ 为奇数；此时 $x_0$ 是\Term{rational-point}{有理点}（2.2.2）。

设 $X_{s(0)}$ 是 $X_s$ 在 $x_0$ 处的\Term{henselization}{Hensel 化}，$\widetilde X_{s(0)}$ 是 $X_{s(0)}$ 在 $x_0$ 处的\Term{blow-up}{爆破}，$Y_0$ 是 $x_0$ 在 $\widetilde X_{s(0)}$ 中的\Term{inverse-image}{逆像}（一个\Term{quadric}{二次超曲面}），且
$X^*_{s(0)}=X_{s(0)}-\{x_0\}=\widetilde X_{s(0)}-Y_0$。
考虑复合\Term{arrow}{箭头}（参见第 XIII 讲 2.4.6.1）：
\[
\begin{tikzcd}[column sep=small]
H^{n-1}(Y_0,\Lambda(m))&H^{n-1}(\widetilde X_{s(0)},\Lambda(m))\arrow[l,"\sim"']\arrow[r]&H^{n-1}(X^*_{s(0)},\Lambda(m))\arrow[d]\\
&H^n_{\{x\}}(X_s,R\Psi_{\bar\eta}(\Lambda(m)))&H^n_{\{x\}}(X_s,\Lambda(m))\arrow[l]
\end{tikzcd}\tag{2.2.6.2}
\]
\begin{SGAresult}{\Term{lemma}{引理}}{2.2.7}
\Term{arrow}{箭头} 2.2.6.2 将 $H^n_{\{x\}}(X_s,R\Psi_{\bar\eta}(\Lambda(m)))$ 等同于 $2m$ 维\Term{quadric}{二次超曲面} $Y_0$ 的 $2m$ 次\Term{cohomology}{上同调}的\Term{primitive-quotient}{本原商}；$\pm\delta$ 是这个\Term{primitive-quotient}{本原商}的自然\Term{generator}{生成元}的像。
\end{SGAresult}
\SourcePage{197}
\Term{cycle}{圈} $\pm\delta$ 由其在 $H_c^i(X_{\bar\eta},\Lambda(m))$ 中的像刻画。利用 (2.1.7.1)，归结为证明复合\Term{arrow}{箭头}
\[
\begin{tikzcd}[column sep=small]
H^{n-1}(Y_0,\Lambda(m))&H^{n-1}(\widetilde X_s,\Lambda(m))\arrow[l,"\sim"']\arrow[r]&H^{n-1}(X_s^*,\Lambda(m))\arrow[d]\\
H_c^n(X_{\bar\eta},\Lambda(m))&H_c^n(X_s,\Lambda(m))\arrow[l,"\operatorname{sp}"']&H^n_{\{x\}}(X_s,\Lambda(m))\arrow[l]
\end{tikzcd}
\]
将第一个群的\Term{primitive-quotient}{本原商}的特定\Term{generator}{生成元}映至最后一个群的特定\Term{generator}{生成元}。

应用 2.1.8，这个\Term{arrow}{箭头}可重写为（精确到\Term{signature}{符号}）
\[
\begin{tikzcd}
H^{n-1}(Y_s,\Lambda(m))\arrow[r,"\partial"]\arrow[d,equal]&H_c^n(X_s,\Lambda(m))\arrow[d]\\
H^{n-1}(Y_{\bar\eta},\Lambda(m))\arrow[r]&H_c^n(X_{\bar\eta},\Lambda(m)),
\end{tikzcd}
\]
由此即得 2.2.7。
\SGAsection{3}{Picard--Lefschetz 公式}
\SGAsection{3.1}{预备结果}
\textbf{3.1.1。} 设 $S$ 是 Hensel \Term{trait}{离散赋值环谱}，$f:X\to S$ 是一个\Term{flat}{平坦}\Term{finite-type}{有限型}\Term{morphism}{态射}，纯\Term{relative-dimension}{相对维数}为 $n$。假设\Term{special-fibre}{特殊纤维} $X_s$ 除\Term{finite}{有限}点集 $\Sigma'$ 外\Term{smooth}{光滑}，在这些点处 $X_s$ 有\Term{ordinary-quadratic-singularity}{普通二次奇点}。
\SourcePage{198}
设 $\Sigma$ 为使\Term{generic-fibre}{一般纤维} $X_\eta$ 在其一个\Term{neighbourhood}{邻域}中\Term{smooth}{光滑}的点 $x\in\Sigma'$ 的\Term{set}{集合}。最后，设 $\Lambda$ 为一个\Term{torsion}{挠性}与 $S$ 的\Term{residue-characteristic}{剩余特征}互素的\Term{torsion-ring}{挠环}。
\begin{SGAresult}{命题}{3.1.2}
(i) 当 $i\ne n$ 时，\Term{sheaf}{层} $R^i\Phi(\Lambda)$ 为零。

(ii) \Term{sheaf}{层} $R^n\Phi(\Lambda)$ 在 $\Sigma$ 之外为零。在 $\Sigma$ 上的\Term{restriction}{限制}是一个\Term{rank}{秩}为 $1$ 的 $\Lambda$-\Term{module}{模}\Term{sheaf}{层}。

(iii) 对 $\Sigma$ 的任意\Term{geometric-point}{几何点} $\bar x$，\Term{form}{形式} $(a,b)$（2.2.5(C)）使 $\Lambda$-\Term{module}{模} $(R^n\Psi(\Lambda))_{\bar x}$ 和 $H^n_{\{\bar x\}}(R\Psi(\Lambda))$ 构成\Term{duality}{对偶}。当 $n\ne0$ 时，它们\Term{free-module}{自由}且\Term{rank}{秩}为 $1$；当 $n=0$ 时，\Term{rank}{秩}为 $2$。
\end{SGAresult}
3.1.2 的断言在\Term{etale-topology}{平展拓扑}下是\Term{local-character}{局部性}的。因此，只须处理 $S$ \Term{strictly-henselian}{严格 Hensel}，且 $X$ 是\Term{relative-affine-space}{相对仿射空间} $\mathbb A_S^{n+1}$ 中由满足 2.2.1(a)(b) 的一个二次方程定义的\Term{hypersurface}{超曲面}的情形（应用 1.3.2）。

$f$ 在 $X_s-\Sigma'$ 的各点\Term{smooth}{光滑}，故在 $\Sigma'$ 之外 $R^i\Phi=0$。由 2.2.4，$R\Phi$ 甚至在 $\Sigma$ 之外为零。现在可假设 2.2.5 的 $(*)$ 成立。断言 (ii) 和 (iii) 于是由 2.2.5(A)(B)(C) 得到。

\SGAsection{3.2}{相对维数为偶数}
沿用 3.1 的记号，取 $n$ 为偶数 $2m$。为简便起见，假设 $S$ 为\Term{strictly-henselian}{严格 Hensel 的}（\Term{generic}{一般}情形可由\Term{descent}{下降}得到）。
\begin{SGAresult}{命题}{3.2.1}
(i) 对 $x\in\Sigma$，群 $H^n_{\{x\}}(R\Psi_{\bar\eta}(\Lambda(m)))$ 有一个精确到\Term{signature}{符号}的自然\Term{generator}{生成元} $\delta$；它由对 $\Lambda$ 的自然性以及下式刻画：
\SourcePage{199}
\[
 (\delta,\delta)=(-1)^m\,2
\]
（当 $n=0$ 时，加上条件 $\operatorname{Tr}(\delta)=0$）。

(ii) 对\Term{inertia-group}{惯性群} $I$ 上某个适当的\Term{character}{特征标} $\epsilon_x:I\to\{\pm1\}$（与 $\Lambda$ 无关），\Term{variation}{变差}写成
\[
 \operatorname{Var}(\sigma)(a)=(-1)^m\frac{\epsilon_x(\sigma)-1}{2}(a\delta)\delta
 \qquad(a\in R^n\Psi_{\bar\eta}(\Lambda(m))).
\]
由此得到
\[
 \sigma(\delta)=\epsilon_x(\sigma)\delta.
\]
\end{SGAresult}
取\Term{limit}{极限}后，只须处理 $\Lambda$ 为 $\mathbb Z_\ell$ 的\Term{versal}{万有}情形。此时满足 (i) 中条件的 $\delta$ 精确到\Term{signature}{符号}至多只有一个。如 3.1.2 中那样，归结到 2.2.5 的 $(*)$ 情形，再应用 2.2.5(D)。

\begin{SGAresult}{补充}{3.2.2}
若对 $x\in\Sigma$，$X$ 在 $x$ 处的\Term{henselization}{Hensel 化} $X_{(x)}$ \Term{isomorphism}{同构}于方程
\[
 \sum_{i\leq j}a_{ij}X_iX_j=0
\]
定义的\Term{projective-scheme}{射影概形} $Q$ 在 $x_0$ 处的\Term{henselization}{Hensel 化}，则\Term{character}{特征标} $\epsilon_x$ 由\Term{fraction-field}{分式域}的\Term{separable-quadratic-extension}{可分二次扩张} $Z(C^+(Q))$ 定义（这由 2.2.5(D) 得到）。
\end{SGAresult}
\textbf{3.2.3。} 当\Term{residue-characteristic}{剩余特征}不为 $2$ 时，可以把 3.2.2 写得更显式。此时 $I$ 有唯一一个非平凡的二阶\Term{character}{特征标} $\epsilon$。对任意\Term{uniformizer}{一致化参数} $t$，有
\[
 \sigma(\sqrt t)=\epsilon(\sigma)\sqrt t.
\]
另一方面，此时 $X_{(x)}$ \Term{isomorphism}{同构}于方程
\SourcePage{200}
\[
 \sum a_{ij}x_i x_j=b
\]
定义的\Term{affine-scheme}{仿射概形}在 $0$ 处的\Term{henselization}{Hensel 化}，其中 $b$ 属于\Term{maximal-ideal}{极大理想}，且 $Q(\underline x)=\sum a_{ij}x_i x_j$ \Term{nondegenerate}{非退化}。\Term{ideal}{理想} $(b)$ 定义 $X_{(x)}$ 中作为 $J^n_{X_{(x)}/S}$ 的\Term{support}{支集}的\Term{subscheme}{子概形}的像。\Term{clifford-algebra}{Clifford 代数}的\Term{centre-algebra}{中心}\Term{isomorphism}{同构}于
\[
 k(\eta)\bigl(((-1)^{m+1}\,2b\det(a_{ij}))^{1/2}\bigr)
\]
（Bourbaki，\textit{Alg.}，第 9 章，§9，第 4 小节末尾的注记），故
\[
 \epsilon_x=\epsilon^{v(b)}.
\]
若 $b$ 的\Term{valuation}{赋值}为偶数，则\Term{variation}{变差}为零；否则 $\epsilon_x=\epsilon$。

\SGAsection{3.3}{相对维数为奇数}
为证明奇\Term{relative-dimension}{相对维数}时的 \Term{picard-lefschetz-formula}{Picard--Lefschetz 公式}，须归结到第 XIV 讲 3.2.1.1 中处理的\Term{transcendental}{超越}情形。所需的\Term{specialization}{特殊化}论证迫使我们在\Term{dimension}{维数}不再为 $1$ 的\Term{base}{基底}上工作。我们将采用针对所需具体情形的专门论证：我手头没有一套在\Term{dimension}{维数}大于 $1$ 的\Term{base}{基底}上可用的理论，也宁愿不用\Term{generic}{一般}性论述来掩盖自己的无知。

\textbf{3.3.1。} 设 $S=\operatorname{Spec}(A)$ 是一个\Term{strictly-local}{严格局部}\Term{noetherian}{Noether}\Term{scheme}{概形}，$f:X\to S$ 是一个\Term{flat}{平坦}\Term{finite-type}{有限型}\Term{morphism}{态射}，纯\Term{relative-dimension}{相对维数}为 $n$，其\Term{special-fibre}{特殊纤维} $X_s$ 的非\Term{smooth}{光滑}点仅为\Term{nondegenerate-quadratic-point}{非退化二次点}。设 $s$ 为 $S$ 的\Term{closed-point}{闭点}，$X_s$ 为\Term{special-fibre}{特殊纤维}，$x_0\in X_s$ 为一个\Term{singular-point}{奇点}。由 1.3.2(i)，在 $x_0$ 附近，$X/S$ \Term{local-scheme}{局部}\Term{isomorphism}{同构}于 $\mathbb A_S^{n+1}$ 中由方程 $Q-b=0$ 定义的一个 $S$-\Term{scheme}{概形}，其中 $b$ 属于\Term{maximal-ideal}{极大理想}，$Q$ 为\Term{nondegenerate}{非退化}\Term{quadratic-form}{二次型}。$A$ 的\Term{ideal}{理想} $(b)$ 只依赖于 $(X/S,x_0)$：在 $x$ 的一个\Term{neighbourhood}{邻域}中，$X$ 在 $\operatorname{Spec}(A/(b))$ 上有一个\Term{section}{截面} $s\in X(A/(b))$，使 $J^n_{X/S}$ \Term{isomorphism}{同构}于 $\mathscr O_{s(\operatorname{Spec}(A/b))}$。将这个\Term{section}{截面}记作 $s_x$，并以 $b(x)$ 表示 $(b)$ 的任意一个\Term{generator}{生成元}。
\SourcePage{201}
\textbf{3.3.2。} 假设 $S$ \Term{normal}{正规}，取 $b\in A$。当 $n$ 在 $A$ 中\Term{invertible}{可逆}且 $b\ne0$ 时，以 $\epsilon_b^{(n)}$ 表示 $A$ 的\Term{fraction-field}{分式域}的\Term{galois-group}{Galois 群}到 $\mathbb Z/(n)(1)$ 的如下\Term{homomorphism}{同态}：
\[
 \epsilon_b^{(n)}(\sigma)=\frac{\sigma\sqrt[n]{b}}{\sqrt[n]{b}}.
\]
有 $(\epsilon_b^{(nm)})^m=\epsilon_b^{(n)}$。若 $\Lambda$ 是一个 $\mathbb Z/n$-\Term{algebra}{代数}，仍以 $\epsilon_b$ 表示复合
\[
 \operatorname{Gal}\longrightarrow\mathbb Z/(n)(1)\longrightarrow\Lambda(1).
\]
它与 $n$ 无关。\Term{character}{特征标} $\epsilon_b$ 通过 $\pi_1(S-(\text{子概形 }b=0))$ \Term{resolution}{分解}。

\textbf{3.3.3。} 假设 $S$ 为\Term{strictly-henselian}{严格 Hensel}\Term{trait}{离散赋值环谱}，$n$ 为奇数：$n=2m+1$。设 $x$ 是 $X_s$ 的一个\Term{singular-point}{奇点}，$X_\eta$ 在它的一个\Term{neighbourhood}{邻域}中\Term{smooth}{光滑}，$X_1/S$ 为 $\mathbb A_S^{n+1}$ 中方程 $Q-b=0$ 定义的一个\Term{subscheme}{子概形}，如 1.3.2(i) 所述。关于 $x$ 的假设意味着：$b$ 模\Term{maximal-ideal}{极大理想}为零，但 $b\ne0$。

设 $X_{s(x)}$ 是 $X_s$ 在 $x$ 处的\Term{henselization}{Hensel 化}，$\widetilde X_{s(x)}$ 是 $\widetilde X_s$ 在 $x$ 处的\Term{blow-up}{爆破}，$Y$ 是 $x$ 的\Term{inverse-image}{逆像}\Term{quadric}{二次超曲面}，且 $X^*_{s(x)}=X_{s(x)}-\{x\}=\widetilde X_{s(x)}-Y$。复合\Term{morphism}{态射}
\[
\begin{tikzcd}[column sep=small]
H^{n-1}(Y,\Lambda(m))&H^{n-1}(\widetilde X_{s(x)},\Lambda(m))\arrow[l,"\sim"']\arrow[r]&H^{n-1}(X^*_{s(x)},\Lambda(m))\arrow[d]\\
&H^n_{\{x\}}(X_s,R\Psi_{\bar\eta}(\Lambda(m)))&H^n_{\{x\}}(X_s,\Lambda(m))\arrow[l]
\end{tikzcd}
\]
将最后一个群等同于\Term{quadric}{二次超曲面} $Y$ 的\Term{cohomology}{上同调}的\Term{primitive-quotient}{本原商}（第 XII 讲 3.5）。事实上，由 (2.2.7)，在一个\Term{local-scheme}{局部}\Term{isomorphism}{同构}的情形中结论成立。这个商的两个自然\Term{generator}{生成元}（$Y$ 的\Term{quadric-generator}{母空间}的\Term{cohomology-class}{上同调类}）的像，是 $H^n_{\{x\}}(X_s,R\Psi_{\bar\eta}(\Lambda(m)))$ 的两个\Term{basis-vector}{基向量} $\pm\delta$。
\SourcePage{202}
\textbf{3.3.4。} 回到 3.3.1 的较\Term{generic}{一般}假设，取 $n=2m+1$。对 $X_s$ 的每个\Term{singular-point}{奇点} $x$，仍以 $\pm\delta_x$ 表示 $R^nf_!\Lambda(m)$ 在 $S$ 上的\Term{section}{截面}，它们是 $H^{n-1}(Y,\Lambda(m))$（记号同上）的\Term{primitive-quotient}{本原商}的\Term{generator}{生成元}在复合\Term{morphism}{态射}
\[
\begin{tikzcd}[column sep=small]
H^{n-1}(Y,\Lambda(m))&H^{n-1}(\widetilde X_{s(x)},\Lambda(m))\arrow[l,"\sim"']\arrow[r]&H^{n-1}(X^*_{s(x)},\Lambda(m))\arrow[d]\\
H^0(S,R^nf_!\Lambda(m))&H_c^n(X_s,\Lambda(m))\arrow[l,"\sim"']&H^n_{\{x\}}(X_s,\Lambda(m))\arrow[l]
\end{tikzcd}
\]
下的像。

$\pm\delta_x$ 是 $x$ 处的\Term{vanishing-cycles}{消失圈}。

当 $S$ 是\Term{trait}{离散赋值环谱}时，$\pm\delta\in H_c^n(X_{\bar\eta},\Lambda(m))$ 是 (3.3.3) 中的\Term{cohomology-class}{类} $\pm\delta\in H^n_{\{x\}}(X_s,R\Psi_{\bar\eta}(\Lambda(m)))$ 在第 XIII 讲 2.1.8.4 下的像。

我们将同时证明下列两个结果。
\begin{SGAresult}{命题}{3.3.5}
若 $S$ 为 Hensel \Term{trait}{离散赋值环谱}，$X/S$ 如 3.3.3 所述，则\Term{variation}{变差}由下式给出：
\[
 \operatorname{Var}(\sigma)(a)=(-1)^{m+1}\epsilon_{b(x)}(\sigma)(a\delta)\delta.
\]
\end{SGAresult}
\begin{SGAresult}{命题}{3.3.6}
设 $S$ 为\Term{strictly-local}{严格局部}\Term{regular}{正则}\Term{scheme}{概形}，$f:X\to S$ 为一个\Term{proper}{固有}\Term{flat}{平坦}\Term{morphism}{态射}，\Term{relative-dimension}{相对维数}为 $n=2m+1$，使得 $X_s$ 恰有一个非\Term{smooth}{光滑}点 $x$，并假设它是\Term{nondegenerate-quadratic-point}{非退化二次点}。再假设 $b(x)$ 是一个参数。则 Galois 在\Term{geometric-generic-fibre}{几何一般纤维}的\Term{cohomology}{上同调}上的\Term{action}{作用}为
\[
 \sigma(a)=a+(-1)^{m+1}\epsilon_{b(x)}(\sigma)(a\delta_x)\delta_x.
\]
\end{SGAresult}
(A) 设 $k$ 是\Term{characteristic}{特征}为 $0$ 的\Term{algebraically-closed}{代数闭}域，$S$ 为 $\operatorname{Spec}(k[T])$ 在 $0$ 处的\Term{henselization}{Hensel 化}。
\SourcePage{203}
当 $b(x)$ 为\Term{uniformizer}{一致化参数}时证明 3.3.5。

由 \Term{lefschetz-principle}{Lefschetz 原理}，可归结到 $k=\mathbb C$。问题是\Term{local-scheme}{局部}的，故可假设 $f:X\to S$ 由 $\mathbb C^{n+1}\to\mathbb C$ 的映射 $\sum z_i^2$ 经\Term{base-change}{基变换}得到。上面给出的 $\pm\delta$ 的构造，在整系数\Term{cohomology}{上同调}中有一个\Term{transcendental}{超越}类比。因此通过第 XIV 讲 2.1，$\pm\delta$ 对应于一个\Term{transcendental}{超越}\Term{vanishing-cycles}{消失圈}，它是
\[
 \pm\delta_{\mathrm{cl}}\in H^n_{\{0\}}(X_0,R\Psi_{\mathrm{cl}\bar\eta}(\mathbb Z))
\]
的像。

由第 XIV 讲 2.1 以及\Term{algebra}{代数}\Term{cycle}{圈} $\pm\delta$ 的相应性质，$\pm\delta_{\mathrm{cl}}$ 在
\[
 H^n_{\{0\}}(X_0,R\Psi_{\mathrm{cl}\bar\eta}(\mathbb Z))\otimes\mathbb Z_\ell
\]
中的像是这个群的一个\Term{generator}{生成元}。因此，$\pm\delta_{\mathrm{cl}}$ 是 $H^n_{\{0\}}(X_0,R\Psi_{\mathrm{cl}\bar\eta}(\mathbb Z))$ 的两个\Term{generator}{生成元}，并与第 XIV 讲 3.2 的\Term{vanishing-cycles}{消失圈}相同。所以公式 3.3.5 由第 XIV 讲 3.2.11 得到（严格说来，应当事先验证第 XIV 讲 2.1 的\Term{isomorphism}{同构}与\Term{cup-product}{杯积}和迹之间的相容性）。

(B) 对 (A) 所述且另外\Term{proper}{固有}的 $X/S$，证明 3.3.6。

这由第 XIII 讲 2.4.6.2 得到。

(C) 设 $p$ 为素数，$S_0$ 为 $\operatorname{Spec}(\mathbb Z[T])$ 在 $(p,T)$ 处的\Term{strict-henselization}{严格 Hensel 化}。对满足 $b(x)=T$ 的 $X/S_0$，证明 3.3.6。

\Term{sheaf}{层} $R^if_*\Lambda$ 在 $S$ 于 $(T)$ 处的\Term{strict-henselization}{严格 Hensel 化}上的\Term{restriction}{限制}适用 (B)。这个\Term{sheaf}{层}在\Term{smooth}{光滑}\Term{divisor}{除子} $T=0$ 的补集上\Term{locally-constant}{局部常值}，且该\Term{divisor}{除子}为\Term{mixed-characteristic}{混合特征}；故 (C) 由 (B) 和 Abhyankar \Term{lemma}{引理}得到。

(D) 在至少一个适用 (C) 的情形中，$\delta_x$ 不是 $\ell$-进\Term{cohomology}{上同调}中的\Term{torsion-element}{挠元素}。
\SourcePage{204}
设 $S_1$ 为 $\operatorname{Spec}(\mathbb Z)$ 在 $(p)$ 处的\Term{strict-henselization}{严格 Hensel 化}，$V\subset\mathbb P^r/S_1$ 为 $S_1$ 上的\Term{smooth}{光滑}\Term{projective-variety}{射影簇}，纯\Term{relative-dimension}{相对维数}为 $n+1$。假设 $V/S_1$ 的\Term{special-fibre}{特殊纤维} $V_p$ 有一个\Term{hyperplane-section}{超平面截面}的\Term{lefschetz-pencil}{Lefschetz 铅笔}（第 XVII 讲 2.5）；这个\Term{lefschetz-pencil}{铅笔}由 $\mathbb P^r/S_1$ 的\Term{dual}{对偶}\Term{projective-space}{射影空间} $\check{\mathbb P}^r/S_1$ 的\Term{special-fibre}{特殊纤维}中的一条直线 $d'$ 定义。将 $d'$ \Term{lifting}{提升}为 $\check{\mathbb P}^r/S_1$ 的一条直线 $d$，设 $X\to d$ 为以 $d$ 参数化的 $V$ 的\Term{hyperplane-section}{超平面截面}\Term{lefschetz-pencil}{铅笔}（第 XVIII 讲 3.1）。$d$ 的例外点集 $\Delta$ 在 $S_1$ 上\Term{etale}{平展}（第 XVII 讲第 6 节）。因此在 $\Delta_p\subset d$ 的一点处作\Term{etale}{平展}\Term{localization}{局部化}，就得到一个 (C) 的情形。

为验证 $\delta_x$ 不是\Term{torsion-element}{挠元素}，只须将标量由 $S_1$ \Term{group-extension}{扩张}到 $\mathbb C$ 后验证。取 $V=\mathbb P^{n+1}$，\Term{projective-embedding}{射影嵌入}由 $\mathscr O(N)$ 定义。当 $N\geq2$ 时，可应用第 XVII 讲 2.5（在这个特殊情形，对所有 $N$ 都成立）。再应用下面的\Term{lemma}{引理}。
\begin{SGAresult}{\Term{lemma}{引理}}{3.3.7}
对 $\mathbb P^{n+1}(\mathbb C)$ 中\Term{degree}{次数} $N$ \Term{large-relative-to-cardinal}{足够大}的\Term{hypersurface}{超曲面}的一个\Term{lefschetz-pencil}{Lefschetz 铅笔}（$N\geq3$ 即足够），一个\Term{vanishing-cycles}{消失圈}（在 $\mathbb Q_\ell$-\Term{cohomology}{上同调}中）非零。
\end{SGAresult}
这个\Term{lemma}{引理}由\Term{lefschetz-pencil}{Lefschetz 铅笔}理论得到，例如由第 XVIII 讲 6.6.1 得到。这里使用第 XVIII 讲这个基于“\Term{picard-lefschetz-formula}{Picard--Lefschetz 公式}”的结果并无循环论证，因为我们已经有 (A)。也可使用第 XIX 讲第 4 小节中更直接的变式，它完全是\Term{transcendental}{超越}的，只依赖于第 XIV 讲 3.2。

(E) 对 $(X/S,x\in X_s)$，在\Term{residue-characteristic}{剩余特征}为 $p$ 时证明 3.3.5。

取 $X_0/S_0$ 如 (D) 所述。由 1.3.2(i)，存在 $g:S\to S_0$，使 $g^*X_0=X_0\times_{S_0}S$ 在 $x_0$ 处的\Term{henselization}{Hensel 化}，与 $X$ 在 $x$ 处的\Term{henselization}{Hensel 化} $S$-\Term{isomorphism}{同构}。因此只须对 $g^*X_0$ 证明 3.3.5。对 $X_0/S_0$ 的 (C) 经\Term{base-change}{基变换}得到
\[
 \sigma(a)=a-(-1)^m\epsilon_{b(x)}(\sigma)(a\delta_x)\delta_x
\]
\SourcePage{205}
在 $H^n(X_{0\bar\eta})$ 中成立，这里 $\sigma$ 属于\Term{inertia-group}{惯性群}。设 $\pi$ 为一个\Term{uniformizer}{一致化参数}，将上式与 (2.2.5.9) 比较，其中 $\epsilon=\epsilon_\pi$，得到
\[
 \lambda_X\epsilon_\pi(\sigma)(a\delta_x)\delta_x
 =-(-1)^m\epsilon_{b(x)}(\sigma)(a\delta_x)\delta_x
\]
在 $H^n(X_{0\bar\eta})$ 中成立。由于 $\delta_x$ 不是\Term{torsion-element}{挠元素}，可在这个公式中约去 $(a\delta_x)\delta_x$；对所得的 $\lambda_X$ 值，公式 2.2.5.9 正是 3.3.5。

(F) 3.3.6 的证明（\Term{generic}{一般}情形）。

由 3.3.5 的\Term{generic}{一般}情形\Term{pushout}{推出} 3.3.6，与由 (A) \Term{pushout}{推出} (B) 和 (C) 相同。

\SGAsection{3.4}{总结}
命题 3.1.2、3.2.1、3.2.2、3.2.3 和 3.3.5 合称为“\Term{picard-lefschetz-formula}{Picard--Lefschetz 公式}”。这个公式是第 XIV 讲 3.2.11 的\Term{algebra}{代数}类比，但根据\Term{relative-dimension}{相对维数}的奇偶性，其\Term{form}{形式}大不相同。

在最值得关注的 $X/S$ \Term{proper}{固有}情形，结果可概括如下。
\begin{SGAresult}{\Term{theorem}{定理}（\Term{picard-lefschetz-formula}{Picard--Lefschetz 公式}）}{3.4}
设 $S$ 为\Term{strictly-henselian}{严格 Hensel}\Term{trait}{离散赋值环谱}，$f:X\to S$ 为一个\Term{proper}{固有}\Term{flat}{平坦}\Term{morphism}{态射}，纯\Term{relative-dimension}{相对维数}为 $n=2m$ 或 $2m+1$。假设 $X_\eta$ \Term{smooth}{光滑}，$X_s$ 只有\Term{ordinary-quadratic-singularity}{普通二次奇点}。设 $\Sigma$ 为 $X_s$ 的\Term{singular-locus}{奇异点集}。

(i) 每个 $x\in\Sigma$ 都附有一个\Term{vanishing-cycles}{消失圈} $\pm\delta_x\in H^n(X_{\bar\eta})(m)$。这些 $\delta_x$ 两两\Term{orthogonal}{正交}。当 $n$ 为偶数时，有 $(\delta,\delta)=(-1)^m\,2$。

(ii) 当 $i\ne n,n+1$ 时，有
\[
 H^i(X_s)\xrightarrow{\sim}H^i(X_{\bar\eta}),
\]
并有\Term{long-exact-sequence}{长正合列}
\SourcePage{206}
\[
\begin{gathered}
0\longrightarrow H^n(X_s)\longrightarrow H^n(X_{\bar\eta})
 \xrightarrow{(\ ,\delta_x)}\bigoplus_{x\in\Sigma}\Lambda(n-m)\\
 \longrightarrow H^{n+1}(X_s)\longrightarrow H^{n+1}(X_{\bar\eta})\longrightarrow0.
\end{gathered}
\]
(iii) 定义了\Term{inertia-group}{惯性群}的\Term{local-scheme}{局部}\Term{character}{特征标}：
\[
\begin{array}{lll}
n\text{ 为偶数}:&\epsilon_x:I\longrightarrow\{\pm1\},\\[2pt]
n\text{ 为奇数}:&\epsilon_x:I\longrightarrow\Lambda(1).
\end{array}
\]
若 $X$ 在 $x$ 处\Term{regular}{正则}，则它们是\Term{surjection}{满射}；若 $x$ \Term{nondegenerate}{非退化}，则它们分别是唯一的二阶\Term{character}{特征标}和\Term{canonical}{典范}\Term{character}{特征标}。

\Term{inertia-group}{惯性群} $I$ 的\Term{action}{作用}由下式给出：
\[
\begin{aligned}
n=2m:\qquad
 \sigma(a)&=a+(-1)^m\sum_{x\in\Sigma}\frac{\epsilon_x(\sigma)-1}{2}(a\delta_x)\delta_x,\\
n=2m+1:\qquad
 \sigma(a)&=a-(-1)^m\sum_{x\in\Sigma}\epsilon_x(\sigma)(a\delta_x)\delta_x.
\end{aligned}
\]
特别地，
\[
\begin{array}{lll}
n\text{ 为偶数}:&\sigma(\delta_x)=\epsilon_x(\sigma)\delta_x,\\[2pt]
n\text{ 为奇数}:&\sigma(\delta_x)=\delta_x.
\end{array}
\]
\end{SGAresult}
